% Binomische Formeln – bank/binomische-formeln/ (zusammenbau v0.4, Heft)
\einheitenkopf[t10]{Binomische Formeln}
\setcounter{aufgabe}{70}

% Anlauf (ohne Prüfkennung)
\begin{aufgabe}{Binomische Formeln – Anlauf}
\begin{teile}
\teil Welche binomische Formel passt, und was ist $a$, was ist $b$ in der Formel? $(x + 14)^2$ \leerfeld Formel; a = \leerfeld , b = \leerfeld
\teil Multipliziere aus: $(x + 4)^2$
\teil Multipliziere aus: $(x - 9)^2$
\end{teile}
\end{aufgabe}

\begin{aufgabe}{Binomische Formeln – Prüfungsaufgaben}
\begin{teile}
\steil Die Parabel $p$ hat die Gleichung $y = (x + 4)^2 - 7$. Weise nach, dass $y = x^2 + 8x + 9$ dieselbe Parabel beschreibt. \hfill (P10 2017 OS)
\steil Lisa behauptet: Der Term $(x - 9)^2 + 2$ lässt sich als $x^2 - 18x + 83$ schreiben. Weise nach, dass sie recht hat. \hfill (P10 2017 OS)
\steil Die Parabel $p(x) = (x - 4)^2 - 3$ und die Gerade $g(x) = -2x + 5$ werden gleichgesetzt. Multipliziere aus und bringe die Gleichung in die Form „… $= 0$“. \hfill (P10 2022 OS)
\steil Die Parabel $p(x) = (x + 1)^2 - 8$ und die Gerade $g(x) = 4x + 8$ werden gleichgesetzt. Multipliziere aus und bringe die Gleichung in die Form „… $= 0$“. \hfill (P10 2022 OS)
\end{teile}
\end{aufgabe}
